% Chapter 6
% Conclusões e Trabalho Futuro

\chapter{Conclusões e Trabalho Futuro}
\label{chap:Chapter6}

Este capítulo sintetiza as conclusões empíricas do Capítulo~\ref{chap:Chapter5}, as limitações e o trabalho futuro.

%----------------------------------------------------------------------------------------

\section{Síntese dos contributos}
\label{sec:conc_contributos}

O trabalho produziu os seguintes contributos: (i) uma definição operacional de \gls{ambivalenciaemocional} a partir da classificação multilabel GoEmotions; (ii) uma \gls{pipeline} documentada e disponibilizada para replicação, em Python, com Hugging Face Transformers, que injecta esse sinal no contexto de geração; (iii) um protocolo cego factorial (condição \(\times\) época) com quatro geradores âncora alinhados a marcos históricos; (iv) um questionário online desenvolvido para este estudo e uma análise centrada no diferencial \(\Delta\) e nas médias absolutas de empatia percebida. As 96 células foram geradas, e a avaliação humana, com 52 participantes no filtro estrito e 2016 classificações na análise principal, está reportada no Capítulo~\ref{chap:Chapter5}.

%----------------------------------------------------------------------------------------

\section{Conclusões}
\label{sec:conc_conclusoes}

Foi construída uma definição operacional de ambivalência a partir do GoEmotions, uma pipeline em Python, com Hugging Face Transformers, que injecta esse sinal no contexto de geração, e um questionário cego com 96 células. Foi testado se esse sinal, face a uma \gls{baseline} sem metadados, aumenta a empatia percebida em quatro geradores âncora.

A hipótese não se confirma. Os quatro contrastes da condição com sinal face à baseline, no modelo misto cruzado, incluem zero no \gls{ic}\,95\% (Tabela~\ref{tab:misto_contrastes}). A variação está entre geradores âncora, cerca de \(1{,}5\) em E1 e cerca de \(6{,}1\) em E3. A pipeline foi implementada e executada com sucesso, permitindo operacionalizar e tornar explícito o sinal de ambivalência; contudo, a sua utilização não produziu melhoria consistente da empatia percebida.

O resultado sustenta apenas que a explicitação de rótulos afectivos derivados do próprio enunciado não apresentou valor incremental detectável neste desenho; não demonstra equivalência entre condições, porque os intervalos de confiança são largos para efeitos pequenos, nem inutilidade da informação afectiva. A interpretação permanece limitada pela ausência de validação humana da ambivalência (Secção~\ref{subsec:disc_constructo}).

O passo seguinte, desenvolvido na Secção~\ref{sec:conc_futuro}, é validar humanamente a detecção de ambivalência, e não repetir o mesmo desenho com mais participantes.

%----------------------------------------------------------------------------------------

\section{Limitações}
\label{sec:conc_limitacoes}

A principal ameaça à validade de constructo decorre da dependência de um único classificador e limiar para definir, seleccionar e representar a ambivalência (Secção~\ref{subsec:disc_constructo}).

As diferenças entre âncoras confundem época, modelo e descodificação; a «época» não foi isolada (Capítulos~\ref{chap:Chapter4} e~\ref{chap:Chapter5}). Uma só resposta por estímulo \(\times\) gerador \(\times\) condição não estima a variabilidade entre outputs do mesmo gerador âncora.

A amostra é de conveniência, por bola de neve na rede pessoal e académica. Iniciaram 156 pessoas; concluíram 63; o filtro estrito retém 52. O abandono é sobretudo precoce. Sem demografia, não se testa se completadores e desistentes diferem. A tradução automática EN\(\rightarrow\)PT foi revista por humano. As respostas dos geradores âncora tiveram apenas uma revisão mínima, alinhada à \gls{qa} automática: células recusadas voltaram a ser executadas na mesma célula até a \gls{qa} aceitar, sem editar o texto nem escolher entre várias válidas. O número de células regeneradas não ficou registado no log. O conteúdo restante não foi editado, para preservar o comportamento gerado.

A cegagem reduz enviesamentos de expectativa, mas não tira a subjectividade da empatia percebida. A opção de itens extra após o mínimo introduz auto-selecção no volume (mediana 24 itens; sete participantes avaliaram as 96 células). O item de prática (\texttt{s01\_\_E1\_\_baseline}) pode reaparecer na rotação pontuada. O ritmo demasiado rápido excluiu quatro participantes no filtro estrito. A verificação de atenção usa um alvo fixo (valor 2). Os testes Welch ao nível da classificação são apenas suplementares; a inferência principal é o misto cruzado (Secção~\ref{subsec:disc_poder}).

A medida de empatia é uma escala global de sete pontos. A escala ordinal foi modelada como contínua na análise principal, por meio de um modelo misto linear, opção comum em escalas de sete pontos com volume elevado de observações; uma replicação poderia verificar robustez com modelos ordinais mistos. Quando as respostas já se situam perto do topo da escala, como em E3 (Secção~\ref{subsec:disc_ancoras}), a sensibilidade a ganhos adicionais diminui. Fala-se aqui em potencial efeito de tecto observado nesta amostra. O desenho avalia respostas isoladas a enunciados isolados. Uma nota elevada num único enunciado não equivale a compreensão social longitudinal \parencite{Zhang2026SAGE}.

A escolha do endpoint, juízo humano cego, e não um modelo a classificar empatia, é uma força do desenho. Comparações recentes entre juízes-\gls{llm} e especialistas encontram inflação nas notas automáticas e menor fiabilidade em dimensões afectivas \parencite{Badawi2026MentalAlign}. Essa força não compensa as ameaças acima; apenas evita uma fonte adicional de viés no critério.

%----------------------------------------------------------------------------------------

\section{Trabalho futuro}
\label{sec:conc_futuro}

O passo que os resultados desta dissertação justificam de forma mais directa é uma validação humana da detecção de ambivalência, independente da geração. Avaliadores, nos mesmos 12 estímulos e antes de qualquer geração, identificam as emoções percebidas, classificam a intensidade de cada uma e avaliam o grau de ambivalência do enunciado. Estima-se a concordância entre avaliadores e entre os juízos humanos e o GoEmotions. Índices como o \(\alpha\) de Krippendorff ou o coeficiente de correlação intraclasse (ICC) ficam como possibilidade, não como método fixo. Sem isso, um \(\Delta\) próximo de zero mistura ``esta operacionalização não teve benefício detectável'' com ``o sinal não mede o constructo''. Qualquer extensão da pipeline deve esperar por essa validação.

Com estímulos assim validados, uma replicação controlada poderia incluir textos não ambivalentes como controlo, várias gerações por célula com parâmetros uniformes, uma escala multi-item de empatia e verificação ordinal de robustez. Extensões como mais âncoras, outras línguas ou variação de \(\tau\) só fazem sentido depois deste controlo.

O resultado observado sugere que o valor, se existir, está menos em explicitar rótulos derivados do próprio enunciado e mais numa representação intermédia (causa, necessidade e estratégia de resposta), a testar só depois da validação humana e da replicação controlada.

Num horizonte posterior, poderá também avaliar-se se histórico conversacional e informação personalizada acrescentam valor que não está disponível no enunciado isolado.

%----------------------------------------------------------------------------------------
